\section{Registration objective}
\label{sec:registration}

Registration (Eq.~\ref{eq:registration}) is solved by differentiable
gradient-based optimisation with PyTorch3D~\cite{pytorch3d}, minimising
$\mathcal{L}=\sum_t\lambda_t\mathcal{L}_t$ over the terms below. The table
states what each term \emph{observes}, which is the information the
optimiser actually has:
\begin{center}
\small
\begin{tabular}{@{}p{0.20\linewidth}p{0.70\linewidth}@{}}
\toprule
\textbf{Term} & \textbf{What it observes} \\
\midrule
Chamfer & Distance from each sampled model point to its nearest target point
  and back; a function of two surfaces, blind to which vertex is which \\
Normal & Agreement of surface normals at matched points \\
Edge / Laplacian / offset & Local mesh regularity and size of residual
  deformation $\mathbf d$; target-independent, ARAP-style~\cite{arap} \\
Symmetry / midline & Bilateral consistency of the model \\
Rotation limits & Joint \emph{angles} against authored ranges, not joint
  positions \\
\bottomrule
\end{tabular}
\end{center}
No term takes a joint location or a vertex identity as an observation.
The surface term is evaluated on $N$ sampled points,
$\mathcal{L}_{\mathrm{ch}}\approx\frac1N\sum_i d(p_i,T)^2$, so the sampling
distribution is itself an implicit weighting over anatomy
(\S\ref{sec:sampling}). Optimisation is staged and hierarchical: global
pose and scale first, then coarse, then fine vertex deformation, rather
than a single flat descent; \S\ref{sec:staged} reports a defect found in this
schedule.
